\section{Conclusion}
\label{sec:conclusion}

Spotlight set out to replace the email-spreadsheet-paper-ballot triplet
that runs most live student showcases with one piece of software that
does simultaneous judging, anonymous voting, integrity moderation, and
real-time admin push; all within a free-tier infrastructure budget
and a single-student maintainer footprint.

The platform meets that goal. Public voters can browse and vote in a
single tap. Judges grade against a segment-specific rubric with auto-save
and an AI-generated reading aid. Admins toggle the live show in real
time and watch event-wide analytics and integrity flags as they happen.
Every privileged write leaves an audit trail; every anonymous vote
passes through five independent integrity defences.

The non-functional posture is strong. Reads sit behind an LRU cache with
explicit invalidation on the writes that matter. Writes are atomic
(Postgres RPC) where they have to be. The error path is uniform:
\texttt{asyncHandler} $\to$ global handler $\to$ structured log line $\to$
JSON response with a request id. SSE plus polling fallback plus an AI
stub plus a degraded health code give the system three lines of
graceful-degradation defence. Feature flags let us disable the riskiest
behaviour without a redeploy.

What remains is an honest debt: no automated test suite, single-region
deploy, process-local fan-out state, no CAPTCHA on the integrity
boundary. None of these compromise correctness today; all of them are
addressable with the iteration paths described in
Section~\ref{sec:future}.

For the contribution as a capstone, the relevant claim is this: a single
small codebase can simultaneously deliver six non-functional qualities
, security, scalability, reliability, maintainability, availability,
efficiency; if the engineering discipline is applied at the
controller boundary (validation, rate-limiting, cache wrapping, audit
logging) rather than retro-fitted later. Spotlight is one demonstration
of that claim.
